\subsection{Model Base Management Subsystem}

This subsystem is the topic assigned to our group, and this section presents it in depth against the full theory: the five elements of the model management subsystem, the functions of the MBMS, the model directory, execution and integration, model selection, and the six required capabilities --- each mapped to our concrete design.

\paragraph{The five elements.}
Turban et al.\ specify the composition of the model management subsystem: it ``\emph{is composed of the following elements}: \emph{model base, MBMS, modeling language, model directory, [and] model execution, integration, and command processor}''~\cite{turban2011}. We implement each element explicitly:

\begin{itemize}
  \item \textbf{Model base} --- the collection of quantitative models that ``\emph{provide the analysis capabilities in a DSS}''~\cite{turban2011}: our five declared models of Section~3.3, plus the \emph{model building blocks and routines} the theory describes --- the regression and smoothing routines of \texttt{statsmodels}, the solver of PuLP, and the optimisation and statistics routines of SciPy --- which are exactly the kind of reusable components (``\emph{random-number generator, curve- or line-fitting routine, present-value computational routine, and regression analysis}'') the theory lists as the contents of a model base.
  \item \textbf{MBMS} --- the management software itself: our registry, described below.
  \item \textbf{Modeling language} --- the means by which models are built; because ``\emph{DSS deals with semistructured or unstructured problems, it is often necessary to customize models, using programming tools and languages}''~\cite{turban2011}, our models are written in Python against a declared interface: each module exposes a run function whose parameter schema is declared in the catalog, so that models can be invoked, combined, and replaced without touching the callers.
  \item \textbf{Model directory} --- ``\emph{a catalog of all the models and other software in the model base. It contains model definitions, and its main function is to answer questions about the availability and capability of the models}''~\cite{turban2011}: this is precisely our \texttt{model\_catalog} table, whose rows carry each model's purpose, variable declaration, parameters, assumptions, and limitations.
  \item \textbf{Model execution, integration, and command processor} --- execution is ``\emph{the process of controlling the actual running of the model}''; integration is ``\emph{combining the operations of several models when needed (e.g., directing the output of one model, say forecasting, to be processed by another one, say a linear programming planning model)}''~\cite{turban2011}: our single entry point \texttt{run\_model(model\_id, params)} controls execution for every caller, and the M1$\to$M2/M3 chaining of Section~3.3 is exactly the integration the theory describes.
\end{itemize}

\paragraph{The four functions of the MBMS.}
The MBMS software ``\emph{has four main functions}: \emph{model creation, using programming languages, DSS tools and/or subroutines, and other building blocks; generation of new routines and reports; model updating and changing; [and] model data manipulation}''~\cite{turban2011}. Our design provides each: \emph{creation} --- a new model joins the suite by writing one Python module against the declared interface and inserting one row into \texttt{model\_catalog}, reusing the building blocks (statsmodels, PuLP); \emph{generation of new routines and reports} --- every scenario run produces a structured result row that the comparison screen turns into reports, and combinations of runs (baseline vs.\ alternative) are new analyses without new code; \emph{updating and changing} --- the catalog row versions a model's declaration and assumptions, so a change is visible, comparable, and revertible; \emph{data manipulation} --- parameters are injected at run time through the registry, which is precisely what makes what-if analysis possible without touching model code.

\paragraph{Model directory details.}
Beyond the definition, the directory answers ``\emph{questions about the availability and capability of the models}''~\cite{turban2011}. Our \texttt{model\_catalog} is queried by the interface to list what models exist, what each answers, and what parameters each accepts; its capability descriptions are what the decision maker reads before choosing a model, and what the run registry validates parameters against.

\paragraph{Model selection.}
Which model to use for which situation ``\emph{cannot be done by the MBMS because it requires expertise and therefore is done manually}''~\cite{turban2011}; our design keeps selection manual --- the revenue manager chooses --- but assists it through the catalog: each row's \emph{purpose} field states the business question the model answers, which is exactly the information a selector needs. The subtler distinction the theory draws --- model selection versus \emph{solution-method selection} (which algorithm solves a chosen model class) --- is handled once at design time: M2 and M4 use closed-form critical ratios because their structure is two-class and single-period, M3 uses an LP solver because the week-level allocation problem has no closed form, and these choices are recorded in the catalog rather than re-decided per run.

\paragraph{Why there is no standard to adopt.}
Turban et al.\ are explicit that ``\emph{model management is an extremely difficult task}'' --- there are hundreds of model classes and many solution algorithms within each class, models are used in different ways in different organisations, and therefore ``\emph{MBMS are generally unique for each type of problem being solved}''; unlike DBMS, there is no standard MBMS~\cite{turban2011}. The corollary for our project is that the model management layer must be designed, not installed --- and its design is the mapping below.

\paragraph{Design: the DBMS analogy, taken literally.}
We take the analogy at the heart of the definition --- an MBMS should relate to models as a DBMS relates to data --- and implement the model base as database tables inside the same DuckDB warehouse that holds the business facts, giving the MBMS a catalog, an execution interface, and a log:

\begin{itemize}
  \item \textbf{\texttt{model\_catalog} --- the directory.} One row per model: identifier, purpose, mathematical type, classification, the four variable roles, parameters, assumptions, limitations, module path, and references. Where a DBMS catalog stores the schema of every table, ours stores the ``\emph{schema}'' of every model.
  \item \textbf{A run registry --- the execution and command processor.} A single entry point, \texttt{run\_model(model\_id, params)}, validates the parameters against the catalog entry, imports the module recorded there, executes the model, and returns structured results. Every caller --- the Streamlit interface, the experiment scripts, the test suite --- goes through the same entry point.
  \item \textbf{\texttt{fact\_scenario\_run} --- the log.} One row per execution: model identifier, scenario name, the full parameter set (JSON), the result (JSON), and the headline KPI --- the memory of every what-if, and the substrate for the comparison reports.
\end{itemize}

A folder of Python modules has none of that: no catalog, no uniform invocation, no run record --- and practice confirms the failure mode, since hotels running revenue management on ``\emph{a series of Excel spreadsheets}'' see success ``\emph{largely dependent on the caliber of the person}'' responsible~\cite{anderson2012}.

\paragraph{The six required capabilities, met one by one.}
The course summarises the expected capabilities of an MBMS as six~\cite{turban2011}; Table~\ref{tab:mbmscaps} maps each to the concrete mechanism in our design.

\begin{table}[H]
\centering
\caption{The six MBMS capabilities and their implementation in this project.}
\label{tab:mbmscaps}
\begin{tabularx}{\textwidth}{@{}l X@{}}
\toprule
\textbf{Capability} & \textbf{How our design provides it} \\
\midrule
Control & All models run exclusively through \texttt{run\_model()}; parameters are validated against the catalog before execution; the model suite is covered by automated tests with known-answer cases (for example, the Littlewood protection level and the overbooking critical ratio on textbook instances). \\
Flexibility & Models compose: M1's forecasts parameterise M2 and M3; any parameter can be overridden at run time to create a what-if scenario; a new model joins the suite by adding a module and a catalog row, without touching the interface or the registry. \\
Feedback & Every run returns its result variables together with the assumptions recorded in its catalog entry; the interface displays the change against a named baseline scenario, so the manager sees not only the number but the direction and size of the move. \\
Interface & The same registry serves the interactive UI, the batch experiment scripts, and the tests --- models contain no interface code, and the interface contains no model logic, so neither duplicates the other. \\
Redundancy reduction & Each model exists exactly once, as one canonical module referenced by its catalog row; analytical SQL over the warehouse is likewise defined once in views that the UI and the models share. \\
Increased consistency & The same model version with the same parameters produces the same results wherever it is run; every result row in \texttt{fact\_scenario\_run} carries its parameter set, so any number in the report can be traced to the run that produced it and reproduced by re-running it. \\
\bottomrule
\end{tabularx}
\end{table}

\paragraph{Why implementing the model base as data matters.}
Three practical consequences follow from treating models as managed data rather than as code. First, \emph{assumptions become visible and centralised}: the walk-cost multiple, the price elasticity, and the capacity estimate are declared once, in the catalog and the dimensions, instead of hiding in source files. Second, \emph{decisions become reproducible}: because each run logs its parameters and results, the comparison screen can answer the retrospective question that Simon's implementation phase poses --- which alternative was chosen, and why. Third, \emph{the reporting stays honest}: a model whose backtest shows it losing to a simple baseline is reported as such, because the run log preserves the comparison rather than the narrative. These are the same reasons organisations built databases in the first place, transposed from data to models.
